\chapter{Tier IV and V: Running Epic Adventures}
łabel{ch:gm−epic}
∈dex{Tier IV}∈dex{Tier V}∈dex{epic tier}∈dex{boons!scarcity at high tiers}

As characters reach Tier IV and V, the game changes. The scope expands from local concerns to regional, national, or world−shaking events. This chapter focuses on how to run these epic tiers effectively, with special attention to how the game \textbf{naturally scales} through Boon scarcity −− not inflated Difficulty Values.

\section{What Changes at Tier IV−−V}

\begin{itemize}
  ℹtem \textbf{Scope:} Actions affect cities, nations, or the entire known world.
  ℹtem \textbf{Consequences:} Failure can shift political borders, end bloodlines, or awaken ancient horrors.
  ℹtem \textbf{Time:} Campaigns often span months or years, with downtime becoming a major phase.
  ℹtem \textbf{Legacy:} Characters become legends −− their deeds echo for generations.
\end{itemize}

The core mechanics remain the same. The difference is \textbf{weight} −− every success and every failure reshapes the fiction in larger, more permanent ways.

\section{Natural Scaling: The Boon Economy Does the Work}

One of Fate's Edge's most elegant features is how it automatically scales difficulty without inflating numbers. At high tiers, characters have large dice pools (often 8−−10 dice). They succeed more often and generate fewer Misses and Partials −− the primary sources of Boons.

\subsection*{Boon Generation Basics (Reminder)}
\begin{itemize}
  ℹtem \textbf{Miss (0 successes):} Gain 2 Boons.
  ℹtem \textbf{Partial (0 < successes < DV):} Gain 1 Boon.
  ℹtem \textbf{Clean Success (successes $≥$ DV, SB = 0):} No Boons.
  ℹtem \textbf{Success with SB (successes $≥$ DV, SB > 0):} No Boons.
\end{itemize}
Boons do not come from winning. They come from \textit{struggling}.

\subsection*{Boon Availability by Tier (Typical)}

\begin{center}
\begin{tabularx}{łinewidth}{lcc}
⊤rule
\textbf{Tier} & \textbf{Avg. Boons earned per scene} & \textbf{Typical use} \\
∣rule
I (0−−40 XP) & 2−−3 & Re−rolls, minor advantages \\
II (41−−90 XP) & 1−−2 & Re−rolls, asset activation \\
III (91−−150 XP) & 1 & Critical rerolls, story shifts \\
IV (151−−220 XP) & 0−−1 & Rare, decisive actions \\
V (221+ XP) & 0−−1 (often 0) & Mythic choices only \\
⊥tomrule
\end{tabularx}
\end{center}

\subsection*{Why This Works}
\begin{itemize}
  ℹtem \textbf{Low−tier characters need the help.} Their small pools produce unreliable results; Boons give them a cushion.
  ℹtem \textbf{High−tier characters must rely on skill, not luck.} With huge pools and reliable success, they no longer have a safety net.
  ℹtem \textbf{The scarcity creates strategic depth.} At high levels, spending a Boon on a re−roll means you cannot activate your spy network later. Every Boon is a meaningful choice.
\end{itemize}

\begin{tcolorbox}[colback=blue!5!white, colframe=blue!75!black, title=Why Brittleness is Good]
∈dex{brittleness}
At high tiers, your power comes from skill, not luck. Fewer Boons mean every roll matters. This is intentional −− embrace the tension.

\medskip
\noindent The simulation insights show that \textbf{chaos players thrived at high tier} because they understood this. They did not hoard Boons; they accepted scarcity and adapted by:
\begin{itemize}
  ℹtem Relying on high Skills and clever positioning.
  ℹtem Using Bonds to generate Boons when needed.
  ℹtem Spending Boons only on critical moments.
\end{itemize}
A \textquotedbl brittle\textquotedbl\ character −− one with few resources but immense competence −− is not a flaw. It is the design. It makes every decision weighty and every success earned.

\medskip
\noindent \textit{Let the dice fall. Let the Boons be precious. Let the tension sing.}
\end{tcolorbox}

\subsection*{Do Not Increase Boon Awards}
Resist the urge to give out extra Boons because players \textquotedbl look cool\textquotedbl\ or \textquotedbl rolled a 10.\textquotedbl\ The scarcity is intentional. It makes high−tier play feel tense and consequential. If you inject Boons, you remove the resource pressure that defines epic play.

\subsection*{Example: Tier I vs. Tier V}
A Tier I rogue picks a lock (DV 3, 5 dice). She rolls a Miss $→$ gains 2 Boons, re−rolls, succeeds.  
A Tier V master thief picks the same lock (DV 3, 10 dice). He rolls a Clean Success $→$ no Boons. Later, he fails a critical roll to disable a ward −− he has only 1 Boon left. He must decide: re−roll this roll, or save it for the escape. The tension is real.

\section{Do Not Inflate Difficulty Values}

\begin{tcolorbox}[colback=red!5!white, colframe=red!75!black, title={Warning: DV Inflation Breaks High−Tier Play}]
∈dex{DV!inflation warning}
A higher DV does not mean \textquotedbl more challenging\textquotedbl\ −− it means \textquotedbl fewer Partials, fewer Boons, faster escalation.\textquotedbl\ At high tiers, players have large pools, but raising DV to 5 or 6 will starve them of Boons entirely and make failure feel arbitrary.

\medskip
\noindent\textbf{Scale Challenge With These Tools Instead:}
\begin{itemize}
  ℹtem \textbf{Timers:} Add a 4−−6 segment timer to the scene (\emph{Palace Coup [4]}, \emph{Ritual Completion [6]}). Players see it tick.
  ℹtem \textbf{Position:} Worsen Position when the fiction demands danger. A Desperate roll already re−rolls successes −− that is plenty.
  ℹtem \textbf{Consequences:} Use SB to add layered costs (see \ref{sec:outcome−sb}).
  ℹtem \textbf{Faction Timers:} Advance long−term threats (\emph{War [8]}) regardless of player success in individual scenes (see \ref{sec:core−resolution−cycle}).
\end{itemize}
\end{tcolorbox}

\subsection*{When Is a Higher DV Appropriate?}
Only when the \emph{fiction} itself demands it: a god's impossible ward, a lock designed by a legendary artisan, a ritual that twists reality. Even then, keep it to DV 5−−6 and telegraph clearly: \textquotedbl This is a mythic challenge. Even with your power, it may fail.\textquotedbl

\section{Failure Must Matter}

At high tiers, a failed roll should never mean \textquotedbl nothing happens.\textquotedbl\ Instead:

\begin{itemize}
  ℹtem \textbf{A rival gains leverage.} \textquotedbl You cannot open the vault, but you see the Duke's sigil on the lock −− he knows someone inside.\textquotedbl
  ℹtem \textbf{A secret becomes public.} \textquotedbl Your investigation fails, but the guard captain mentions your name to a reporter.\textquotedbl
  ℹtem \textbf{A faction timer advances.} \textquotedbl The ritual continues. \emph{Summoning [8]} ticks +1.\textquotedbl
  ℹtem \textbf{An asset is compromised.} \textquotedbl Your safehouse is discovered −− it is now \emph{Neglected}.\textquotedbl
\end{itemize}

If failure merely delays action without changing the situation, you are not using the system to its full potential.

\section{Using Bonds as the Primary Boon Engine}

At high tiers, Boons should mainly come from \textbf{Bond−driven actions} (see \ref{subsec:boons−earn}). Encourage players to:
\begin{itemize}
  ℹtem Call on Bonds in difficult moments: \textquotedbl I owe you −− I ask for your help now.\textquotedbl
  ℹtem Expose allies or followers to danger to gain Boons.
  ℹtem Accept new obligations or narrative fallout in exchange for Boons.
\end{itemize}

\subsection*{Example: Bonded Allies Transferring Boons}
The party faces a mythic guardian. Sariel, the rogue, has 0 Boons left. Kestra, the mage, has 2 Boons but is pinned at range. They share a Bond (\textquotedbl She saved me from the shadow cult.\textquotedbl)

\begin{quote}
  Kestra: \textquotedbl I spend my action to transfer 2 Boons to Sariel. As I channel, I shout: `Remember the Forge −− you owe me!'\textquotedbl\\
  GM: \textquotedbl The Bond flares. Sariel receives 2 Boons. Spend them wisely.\textquotedbl
\end{quote}
Sariel uses 1 Boon to re−roll a failed stealth check and 1 Boon to improve her Position to Dominant. She slips past the guardian and disables the trap.

This turns a Bond from backstory into a tactical resource. At high tiers, such exchanges are the difference between victory and defeat.

\section{Designing High−Tier Scenes}

When planning an epic scene, ask:
\begin{enumerate}
  ℹtem \textbf{What changes if the characters succeed?}
  ℹtem \textbf{What compromises emerge on a Partial success?}
  ℹtem \textbf{What new dangers arise on a Miss?}
\end{enumerate}

If the answer to \textquotedbl Miss\textquotedbl\ is \textquotedbl they try again,\textquotedbl\ increase the stakes, not the DV.

\subsection*{Example Scene: The Royal Vote}
The party tries to convince the council to declare war on a neighboring realm.
\begin{itemize}
  ℹtem \textbf{Clean Success:} The vote passes; the army marches under their banner.
  ℹtem \textbf{Success with SB:} The vote passes, but a rival is appointed general.
  ℹtem \textbf{Partial:} The vote is delayed; new timer \emph{Council Deadlock [4]} appears.
  ℹtem \textbf{Miss:} The vote fails. The rival declares war anyway −− without the party's leadership.
\end{itemize}
No DV was inflated. The stakes did the work.

\section{Pacing and Scope}

\subsection*{Embrace Longer Timelines}
Time becomes a tool. Use downtime to:
\begin{itemize}
  ℹtem Advance faction timers.
  ℹtem Allow rumors to spread.
  ℹtem Let players build legacies (fortresses, guilds, schools).
\end{itemize}

\subsection*{Abstract the Mundane}
At high tiers, do not roll for climbing walls or picking common locks. Assume competence. Save rolls for choices that matter.

\subsection*{Let Players Command Armies}
Treat armies as high−Cap followers (see \ref{sec:followers−assets−gm}). Use \textbf{Battle Timers} (4−−6 segments) for skirmishes. Do not resolve each soldier's attack −− abstract with leadership rolls.

\section{Caveats and Warnings}

\begin{itemize}
  ℹtem \textbf{Do not design adventures like Tier I.} Linear dungeons and small−scale fetch quests feel beneath legendary heroes. Instead, give them systemic problems: a curse that erases memory, a political schism, a dying god.
  ℹtem \textbf{Do not hoard Boons.} Spend SB and optional Threat Pool points (see) to create memorable moments. Players earned their power −− challenge them with creative complications, not starvation.
  ℹtem \textbf{Do not fall into \textquotedbl superhero\textquotedbl\ mode.} The world should still react believably. A legendary warrior can be assassinated in their sleep. A master mage can be outmaneuvered politically.
\end{itemize}

\section{Tools for Managing Boon Scarcity}

High−tier play demands that GMs and players rethink how Boons are used. The following tools help make scarcity feel tense, not frustrating.

\subsection*{Boon−Worthy Moments (Player Guidance)}
At Tier IV−−V, a Boon should never be spent on a trivial re−roll. Instead, reserve Boons for:
\begin{itemize}
  ℹtem \textbf{Negating a catastrophic failure} (e.g., re−rolling a Miss that would kill an ally).
  ℹtem \textbf{Activating a key Asset} (e.g., calling in a spy network to reveal a weakness).
  ℹtem \textbf{Shifting Position from Desperate to Controlled} before a mythic roll.
  ℹtem \textbf{Powering a Prestige or Epic talent} that changes the scene.
\end{itemize}
\textbf{GM tip:} Remind players of this at the start of each high−tier session.

\subsection*{The \textquotedbl Substituting Fatigue for Boons\textquotedbl\ Rule (Revisited)}
Recall that a player may pay 1 Fatigue instead of 1 Boon (see Section \ref{subsec:boons−spend}). At high tiers, this becomes a lifeline. Encourage players to use it, but attach a narrative cost: \textquotedbl You push through the pain, but your hand trembles −− your next roll is at Desperate Position.\textquotedbl

\subsection*{Follower Sacrifice as Boon Substitute}
A loyal follower (Cap 3+) can take the consequence of a failed roll in place of a Boon. The follower suffers Harm 1 (or a Condition) and the player avoids spending a Boon. This should be used sparingly (once per session maximum) and only when the fiction supports it.

\subsection*{Let the Roll Stand}
At high tiers, players may reflexively reach for Boons to re−roll one failure. Sometimes, the most dramatic outcome is to \textbf{let the failure stand}. As GM, you can say: \textquotedbl You have 0 Boons left. The die is cast. What does failure look like?\textquotedbl\ This forces players to engage with consequences rather than erase them.

\subsection*{Example: A High−Tier Turn with 0 Boons}
The party's archmage, Jora (Tier V), attempts to counterspell a world−ending ritual. She has 0 Boons. She rolls 8 dice (Wits 5 + Arcana 3) vs DV 4. Result: 3 successes −− a Partial.
\begin{itemize}
  ℹtem \textbf{Without Boons:} She cannot re−roll. The GM rules: \textquotedbl The ritual slows, but a backlash rips through the ley lines. A nearby village is consumed by wild magic. You save the world, but at a cost.\textquotedbl
  ℹtem \textbf{If she had a Boon:} She could re−roll a failure and possibly achieve a Clean Success, avoiding the backlash.
\end{itemize}
The scarcity of Boons made the Partial outcome feel meaningful, not like a failure.

\subsection*{Physical Trackers for Boon Scarcity}
To make Boon scarcity visible and tactile:
\begin{itemize}
  ℹtem Use a \textbf{small bowl or tray} in the center of the table with poker chips representing the \emph{party's} total Boons (players still track individual Boons, but a shared visual reinforces scarcity).
  ℹtem When a player spends a Boon, they physically move a chip to a \textquotedbl spent\textquotedbl\ pile. When they earn a Boon, they add a chip.
  ℹtem For high−tier sessions, set a \textbf{table limit of 3 Boons maximum per player} (down from 5). This makes each chip precious.
\end{itemize}

\section{Optional Epic Tools (Brief)}

The following are covered in other chapters or appendixes:

\begin{itemize}
  ℹtem \textbf{Boss Generator} −− See Appendix F (in the full GM's Guide) for a card−based tool to create mythic adversaries.
  ℹtem \textbf{Legacy Projects} −− See Chapter \ref{ch:gm−resources} for long−term timers that span sessions.
  ℹtem \textbf{Mythic Martial Tags} −− Player−facing epic talents for warriors; see Chapter 6 of the \textit{Player's Guide}.
  ℹtem \textbf{Mass Combat} −− See optional rules in Appendix G for battle timers and army followers.
\end{itemize}

\begin{tcolorbox}[colback=red!5!white, colframe=red!75!black, title={The Mythic Failure Table (d12)}]
∈dex{mythic failure}

When a high−tier character fails a pivotal roll, describe a consequence from this table (or invent your own).

\small
\begin{tabularx}{łinewidth}{lX}
⊤rule
\textbf{d12} & \textbf{Consequence} \\
∣rule
1 & An innocent dies because of your failure. \\
2 & A rival claims the victory instead of you. \\
3 & An artifact shatters, releasing a curse or a new threat. \\
4 & Your reputation is damaged −− lose 1 Mandate (if using, see \ref{subsec:mandate}). \\
5 & An ally is captured or turned against you. \\
6 & The problem spreads −− a new timer appears (\emph{Crisis [4]}). \\
\addlinespace
7 & A Patron withdraws favor −− lose access to one Rite until you atone. \\
8 & A faction declares you an enemy −− start a \emph{Vendetta [6]} timer. \\
9 & The land itself is wounded −− a region becomes blighted or uninhabitable. \\
10 & You lose a memory or a cherished belonging (GM chooses). \\
11 & A prophecy twists −− your success now serves a dark purpose you cannot yet see. \\
12 & Your failure opens a door −− something ancient and hungry awakens. \\
⊥tomrule
\end{tabularx}
\normalsize
\end{tcolorbox}

\section{Sample High−Tier Scenario: The Dying Sun}

\textbf{Premise:} The realm's patron sun−spirit is fading. Without it, crops fail, undead rise, and the land dies.

\textbf{Scope:} Continental.

\textbf{Timers:}
\begin{itemize}
  ℹtem \emph{Sun's Fading [10]} −− ticks each session, faster if neglected.
  ℹtem \emph{Cult of Eternal Night [6]} −− a faction that worships the darkness.
  ℹtem \emph{Alliance of the Dawn [8]} −− political support for a ritual to save the sun.
\end{itemize}

\textbf{Success:} The party finds a lost relic, performs a ritual, and restores the sun −− but must choose a cost (e.g., their own lifespans, a Patron's demand).

\textbf{Partial:} The sun is saved, but the cult is not destroyed −− they will return.

\textbf{Miss:} The sun dies. The party must lead a desperate exodus or accept a dark bargain for a new light.

This scenario cannot be solved by a single roll. It requires investigation, diplomacy, combat, and sacrifice −− perfect for epic tier play.

\section{Final Thought}

At Tier IV and V, your players are no longer adventurers −− they are \textbf{legends in the making}. Honor their power by giving them problems that require wisdom, sacrifice, and creativity to solve.

The dice still matter. Consequences still flow. But now, every choice echoes across nations and generations.

Make it legendary.